\documentclass[9pt]{article}
\usepackage[a4paper,margin=0.25in]{geometry}
\usepackage{amsmath,amssymb,amsfonts,mathtools}
\setlength{\jot}{4pt} 
\usepackage{braket}
\usepackage{slashed}
\usepackage{cancel}


\usepackage{graphicx}
\usepackage{tikz}
\usepackage{pgfplots}


\usepackage{xcolor}
\usepackage{hyperref}
\usepackage{comment}


\newcommand{\CGCket}{|{\rm CGC}\rangle}
\newcommand{\CGCbra}{\langle {\rm CGC}|}

\def\wt#1{\widetilde{#1}}
\def\mf#1{\mathbb{#1}}

\newcommand{\A}[3]{\mf{A}^{#2}_{#3}(#1)}

\numberwithin{equation}{section}

\title{Terms proportional to scales}
\author{}
\date{}

\begin{document}
\maketitle {\color{red}[Marks of the check: red = corrected or added (part of the expression), red text in square brackets = notes, }{\color{teal}\cancel{\text{teal, struck}}}{\color{red}\ = removed. Green, orange and blue are the author\textquotesingle s own highlighting and are unchanged.]}
\section{Group I (with c.c.)}
\subsection{Row I}
\textsc{NA}
\subsection{Row II}
\textsc{NA}
\subsection{Row III}
\begin{align}
   \frac{d^{3}N_{1\,\rm NLO}^{\rm}}{d^{3}k}  & = \frac{1}{(2\pi)^{3}}  \int_{w, w^\prime} e^{-ik{\color{teal}\cancel{(w^\prime-w)}}{\color{red}\cdot(w-w^\prime)}}\, \Bigg\{\notag \\ & \frac{ig^{4}\,f^{abd}}{8\pi^{5}}\,\int_{x,y,z,x'}\, \frac{(x'-w')\cdot (y-z)}{(x'-w')^{2}\,(y-z)^{2}}\, \frac{1}{\left[(x-z)^{2}-(x-w)^{2}\right]}\times\, {\color{green}\left[\frac{1}{\Lambda}+\frac{1}{k^{+}-\vee}\right]}\notag \\ & \left[U^{ab'}(x')\,\rho^{b'}(x')-U^{ab'}(w')\, \rho^{b'}(x')\right]\times\, \left[U^{de}(x)\,U^{bc}(z)\, \rho^{e}(x)\,\rho^{c}(y)-U^{de}(x)\,U^{bc}(y)\,\rho^{e}(x)\,\rho^{c}(y)\right]\notag \\ &  + \frac{ig^{4}\,f^{abd}}{8\pi^{5}}\,\int_{x,y,z,x'}\, \frac{(x'-w')\cdot (y-z)}{(x'-w')^{2}\,(y-z)^{2}}\, \frac{1}{\left[(x-z)^{2}-(x-w)^{2}\right]}\times\, {\color{green}\left[\frac{1}{\Lambda}-\frac{1}{k^{+}-\Lambda}\right]}\notag \\ & \left[U^{ab'}(x')\,\rho^{b'}(x')-U^{ab'}(w')\, \rho^{b'}(x')\right]\times\, \left[U^{de}(x)\,U^{bc}(z)\, \rho^{e}(x)\,\rho^{c}(y)-U^{de}(x)\,U^{bc}(y)\,\rho^{e}(x)\,\rho^{c}(y)\right]\Bigg\}
\end{align} {\color{red}[Checked: both green brackets are the double-pole parts of the first two terms of the $p^{+}$ integrand of Row III: with $D=p^{+}(x-z)^{2}-k^{+}(x-w)^{2}$, $\frac{p^{+}+k^{+}}{(k^{+}-p^{+})^{2}\,D}$ $=\frac{2}{(x-z)^{2}-(x-w)^{2}}\,\frac{1}{(p^{+}-k^{+})^{2}}$ $+$ simple poles (exact partial fractions), $\int_{k^{+}+\Lambda}^{\vee}\frac{dp^{+}}{(p^{+}-k^{+})^{2}}=\frac{1}{\Lambda}-\frac{1}{\vee-k^{+}}$, $\int_{\Lambda}^{k^{+}-\Lambda}\frac{dp^{+}}{(p^{+}-k^{+})^{2}}=\frac{1}{\Lambda}-\frac{1}{k^{+}-\Lambda}$. Agrees with $N_{1}$ of the checked Group I file. Phase: convention $e^{-ik\cdot(w-w^\prime)}$ (here and below).]}\par\noindent
\begin{align}
   \frac{d^{3}N_{2\,\rm NLO}^{\rm}}{d^{3}k}  & = -\,\frac{1}{(2\pi)^{3}}  \int_{w, w^\prime} e^{-ik{\color{teal}\cancel{(w^\prime-w)}}{\color{red}\cdot(w-w^\prime)}}\, \Bigg\{\notag \\ & \frac{ig^{4}\, f^{abd}}{8\pi^{5}}\, \int_{x,y,z,x'}\,\frac{(x'-w')\cdot (y-z)}{(x'-w')^{2}\, (y-z)^{2}\,\left[(x-w)^{2}-(x-z)^{2}\right]}\,{\color{green}\left[\frac{1}{(\vee-k^{+})}-\frac{1}{\Lambda}\right]} \notag \\ & \left[U^{ab'}(x')\,\rho^{b'}(x')-U^{ab'}(w')\, \rho^{b'}(x')\right]\times\,\left[U^{de}(x)\,U^{bc}(z)\, \rho^{e}(x)\,\rho^{c}(y)-U^{de}(x)\,U^{bc}(y)\,\rho^{e}(x)\,\rho^{c}(y)\right]\Bigg\} 
\end{align} {\color{red}[Checked: the $p^{+}>k^{+}$ non-instantaneous term has the double pole $-\frac{1}{(x-z)^{2}-(x-w)^{2}}\frac{1}{(p^{+}-k^{+})^{2}}$ (only the $\delta_{ij}\delta_{km}$ structure contributes; $(z-w)\cdot[(x-w)-(x-z)]=(z-w)^{2}$), range $k^{+}+\Lambda<p^{+}<\vee$; with the overall minus sign this is the expression above.]}\par\noindent
\begin{align}
\frac{d^{3}N_{3a\,\rm NLO}}{d^{3}k} & =\frac{1}{(2\pi)^{3}}\int_{w,w'}e^{-ik{\color{teal}\cancel{(w'-w)}}{\color{red}\cdot(w-w')}}\Bigg\{\notag \\ & \frac{ig^{4}f^{abd}}{8\pi^{5}}\,{\color{green}\left[\frac{1}{\Lambda}-\frac{1}{k^{+}-\Lambda}\right]}\,\int_{x,y,z,x'}
\frac{(x'-w')\cdot(y-z)}{(x'-w')^{2}(y-z)^{2}\,\left[(x-w)^{2}-(x-z)^{2}\right]}\notag \\ & \left[U^{ab'}(x')\,\rho^{b'}(x')-U^{ab'}(w')\, \rho^{b'}(x')\right]\times \left[U^{de}(x)\,U^{bc}(z)\, \rho^{e}(x)\,\rho^{c}(y)-U^{de}(x)\,U^{bc}(y)\,\rho^{e}(x)\,\rho^{c}(y)\right]\Bigg\}
\end{align} {\color{red}[Checked: the $p^{+}<k^{+}$ non-instantaneous term has the same double pole $-\frac{1}{(x-z)^{2}-(x-w)^{2}}\frac{1}{(p^{+}-k^{+})^{2}}$, range $\Lambda<p^{+}<k^{+}-\Lambda$. $N_{3b}$ and the instantaneous (third) term have only simple poles: no power terms.]}\par\noindent
\subsection{Row IV}
\textsc{NA}
\subsection{Row V}
\textsc{NA} {\color{red}[No inverse power of a cut-off. For completeness: $N_{2b\,\rm UV}$ has the coefficient $\frac{1}{k^{+}}\big[2\log\frac{k^{+}-\Lambda}{\Lambda}-\frac{11}{6}+\frac{4\Lambda}{k^{+}}-\frac{\Lambda^{2}}{k^{+2}}+\frac{2\Lambda^{3}}{3k^{+3}}\big]$: positive powers of $\Lambda$, which vanish for $\Lambda\to0$.]}
\section{Group II}
\subsection{Row I}
\textsc{NA} {\color{red}[No $1/\Lambda$. But $N_{2f}$ contains a pure power without logarithm, $\frac{1}{(2\pi)^{3}}\frac{g^{4}d_{\perp}}{32\pi^{5}}\big[\frac{1}{\vee}-\frac{1}{k^{+}+\Lambda}\big]\int\frac{e^{-ik\cdot(w-w^\prime)}}{(z-w)^{2}(z-w^\prime)^{2}}[\dots]$ (from the double pole at $p^{+}=-k^{+}$, $\int_{\Lambda}^{\vee-k^{+}}\frac{dp^{+}}{(p^{+}+k^{+})^{2}}$); it tends to $-\frac{1}{k^{+}}$, a finite term, and the $\frac{1}{\vee}$ part vanishes. To be listed here or not: for the author.]}
\subsection{Row II}
\textsc{Only 4 $\rho$'s}
\subsection{Row III}
\textsc{Only 4 $\rho$'s}
\subsection{Row IV}
\textsc{NA} {\color{red}[No $1/\Lambda$. The BFKL coefficient is $\frac{1}{k^{+}}\log\frac{\vee}{k^{+}}-\frac{1}{k^{+}}+\frac{1}{\vee}$ ($=\frac{1}{k^{+}}\int_{k^{+}/\vee}^{1}d\xi\,\frac{1-\xi}{\xi}$): the $\frac{1}{\vee}$ vanishes for $\vee\to\infty$ and $-\frac{1}{k^{+}}$ is a finite term. To be listed here or not: for the author.]}
\section{Group III\, (with c.c.)}
\subsection{Row I}
\textsc{NA}
\subsection{Row II}
\textsc{NA}
 \subsection{Row III}
 \textsc{NA}
\subsection{Row IV}
\textsc{NA}
\subsection{Row V}
\textsc{NA}
\subsection{Row VI}
\textsc{NA}
\section{Algebra}
\begin{align}
   \frac{d^{3}N_{\rm NLO}^{\rm}}{d^{3}k}  & = \frac{1}{(2\pi)^{3}}\,\frac{ig^{4}\,f^{abd}}{8\pi^{5}}\,{\color{red}\int_{w,w'}}\,\int_{x,y,z,x'}\, e^{-ik{\color{teal}\cancel{(w^\prime-w)}}{\color{red}\cdot(w-w^\prime)}}\,\frac{(x'-w')\cdot (y-z)}{(x'-w')^{2}\,(y-z)^{2}}\, \frac{1}{\left[(x-z)^{2}-(x-w)^{2}\right]}\times\,\notag \\ &  \left[{\color{orange}\frac{1}{\Lambda}}+\frac{1}{k^{+}-\vee}+{\color{orange}\frac{1}{\Lambda}}{\color{blue}-\frac{1}{k^{+}-\Lambda}}+\frac{1}{\vee - k^{+}}{\color{orange}-\frac{1}{\Lambda}-\frac{1}{\Lambda}}{\color{blue}+\frac{1}{k^{+}-\Lambda}}\right]\times\notag \\ & \left[U^{ab'}(x')\,\rho^{b'}(x')-U^{ab'}(w')\, \rho^{b'}(x')\right]\times\, \left[U^{de}(x)\,U^{bc}(z)\, \rho^{e}(x)\,\rho^{c}(y)-U^{de}(x)\,U^{bc}(y)\,\rho^{e}(x)\,\rho^{c}(y)\right]  = 0
\end{align} {\color{red}[Checked: the left-hand side is the sum of the power terms of $N_{1}+N_{2}+N_{3a}$. With the common denominator $(x-z)^{2}-(x-w)^{2}$ the factors are $N_{1}$: $\frac{1}{\Lambda}+\frac{1}{k^{+}-\vee}+\frac{1}{\Lambda}-\frac{1}{k^{+}-\Lambda}$; $N_{2}$: $\frac{1}{\vee-k^{+}}-\frac{1}{\Lambda}$; $N_{3a}$: $-\frac{1}{\Lambda}+\frac{1}{k^{+}-\Lambda}$. The sum of the eight is identically zero at finite $\Lambda$ and $\vee$ (not only in the limit). The integral over $w,w^\prime$ was missing in front of the phase.]}\par\noindent
\end{document}